
One generation IV reactor design is the \gls{MSR} \cite{kelly_generation_2014}.
\Glspl{MSR} are advanced reactors intended to improve safety and efficiency over traditional reactor designs.
Instead of a solid fuel, \glspl{MSR} use a liquid fuel composed of a molten salt with some fissile nuclide, such as $^{235}$U.
\Glspl{MSR} offer many benefits: high operating temperature, low operating pressure, and high thermal conductivity and capacity of the salts.
Other benefits include a high coefficient of thermal expansion for negative temperature coefficient of reactivity, high resource utilization through higher burn-up, reduced cost of fabricating and transporting fuel elements, and improved neutron economy by online fission product removal \cite{serp_molten_2014}.

\Glspl{MSR} have several barriers to deployment which must be addressed before they are commercialized.
Corrosion and material degradation pose an issue, as the high temperature salts corrode materials.
Another deployment barrier involves reactor safety, security, and safeguards by design, as \glspl{MSR} have distinct designs compared with traditional light water reactors.
To properly design \glspl{MSR}, models must be able to simulate them accurately.
One specific phenomena of critical importance in \glspl{MSR} are the circulating \glspl{DNP}.

\Glspl{DNP} are nuclides which undergo $\beta^-$ decay, producing a daughter, called an emitter, which emits a neutron \cite{bank2000jef}.
Without properly modeling the \glspl{DNP} in reactors, simulations may predict incorrect power responses, reactivity worth, and reactor controllability \cite{das1993importance, das1977kinetics, lane1958fluid}.
Simulations typically do not model \glspl{DNP} explicitly, but instead represent them using \gls{DNP} group parameters, which preserve the same physical effect at much lower computational cost.
Tables \ref{tab:dnp-timeline-1} and \ref{tab:dnp-timeline-2} show a brief historical timeline of \gls{DNP} group parameter modeling advancements.


\begin{table}[!ht]
\caption{Timeline of advancements in \gls{DNP} modeling up to 1982.}
% https://tex.stackexchange.com/questions/196794/how-can-you-create-a-vertical-timeline
\centering
\begin{minipage}[t]{.7\linewidth}
\color{gray}
\rule{\linewidth}{1pt}
\ytl{1939}{Roberts et al. discover delayed neutrons \cite{roberts1939further}; Bohr and Wheeler explain delayed neutrons \cite{bohr1939mechanism}}
\ytl{1946}{Burgy et al. conduct initial spectral measurements \cite{burgy1946energies}}
\ytl{1947}{Hughes et al. initial use of six \gls{DNP} groups at Argonne heavy water pile \cite{hughes1948delayed}}
\ytl{1957}{Keepin et al. generate improved \gls{DNP} group parameters \cite{keepin1957delayed, williams1996choice}}
\ytl{1965}{Amiel identifies seven \glspl{DNP} that account for 80\% of delayed neutrons \cite{amiel1965delayed}}
\ytl{1969}{Marmol compiles experimental data on emission probability and half-life for 34 \glspl{DNP} \cite{del1969delayed}}
\ytl{1972}{Tomlinson maps 38 \glspl{DNP} onto six groups \cite{tomlinson1972theory}}
\ytl{1977}{Saphier et al. construct delayed neutron spectra from 20 \glspl{DNP} \cite{saphier1977evaluated}}
\ytl{1982}{Rudstam generate \gls{DNP} group parameters and spectra using 67 \glspl{DNP} based on half-lives \cite{rudstam1982six}}
\bigskip
\rule{\linewidth}{1pt}%
\end{minipage}%
\label{tab:dnp-timeline-1}
\end{table}


\begin{table}[!ht]
\caption{Timeline of advancements in \gls{DNP} modeling from 1983 onward.}
% https://tex.stackexchange.com/questions/196794/how-can-you-create-a-vertical-timeline
\centering
\begin{minipage}[t]{.7\linewidth}
\color{gray}
\rule{\linewidth}{1pt}
\ytl{1983}{England et al. generate \gls{DNP} group parameters and spectra using 105 \glspl{DNP} based on half-lives \cite{england1983aggregate}}
\ytl{1988}{Brady generates \gls{DNP} group parameters and spectra using 271 \glspl{DNP} using least squares \cite{brady1988evaluation}. Group parameters in ENDF/B-VI.8 set and remain unchanged through ENDF/B-VIII.0 \cite{wilson2002delayed, Brown20181, endf6, brady1989delayed, england1988status, brady1989validation}}
\ytl{1999}{Parish et al. state \gls{DNP} group parameters in \gls{ENDF} are in need of revision \cite{parish1999status}}
\ytl{2014}{\gls{IAEA} \gls{CRP} on Development of a Reference Database for
$\beta$-delayed neutron emission data \cite{mills2014calculation}}
\ytl{2020}{Leconte et al. use more than 400 \glspl{DNP} to combine data with macroscopic experimental measurements \cite{leconte2020new}}
\ytl{2025}{This study develops \gls{DNP} group parameters for use in \glspl{MSR}}
\bigskip
\rule{\linewidth}{1pt}%
\end{minipage}%
\label{tab:dnp-timeline-2}
\end{table}

\gls{DNP} modeling and group parameter formulation advanced rapidly thanks to increased experimental data for individual \glspl{DNP}, which enabled more accurate group parameter formulation.
However, these formulations of \gls{DNP} group parameters do not account for certain physics unique to \glspl{MSR}.
Specifically of interest is how the fuel movement and chemical reprocessing affects the formulation of \gls{DNP} group parameters.
Although researchers have analyzed \glspl{DNP} in \glspl{MSR}, that work has not investigated this gap \cite{zhang2009analysis, zuo2022flow, lee2022transient, shi2021gen, kupriyanov2022effective}.
Most of these studies investigate \glspl{DNP} movement during transients.
However, accurate transient models require accurate \gls{DNP} group parameters.
Thus, the next step in group parameter development is \gls{DNP} group parameters that accurately reflect physical phenomena in \glspl{MSR}.

This research aims to generate improved \gls{DNP} group parameters that account for \gls{MSR} physical phenomena, verify the parameters with the literature, validate the parameters by simulating transients and comparing with experimental data, and evaluate the impact of the improved parameters on safety, safeguards, and security by design.
For example, chemical removal of \glspl{DNP} decreases the number of delayed neutrons emitted per fission event on average.
This reduces the total delayed neutron yield, decreases the effective delayed neutron fraction, shortens the reactor period, and reduces reactor controllability.
Current \gls{DNP} group parameters in \glspl{MSR} simulations may overestimate the number of delayed neutrons in the system.

The physical phenomena in \glspl{MSR} can be incorporated at varying levels of fidelity, with associated differences in cost and accuracy.
This document describes three models: a 0D scaled model, already complete; a 0D flow model, proposed; and a spatially resolved model, discussed for future work.
The least accurate and fastest model is the 0D scaled model.
This model neglects time dependence, decay chains, and cross-sections while scaling the flux and chemical removal rates based on the fraction of fuel salt in different regions of the primary loop.
The intermediate approach uses a 0D flow model.
This model accounts for time dependence and decay chains while applying the flux and chemical removal rates only in the relevant parts of the primary loop.
The most accurate and slowest approach uses a spatially resolved model.
This model accounts for time and space dependence across the entire primary loop.
Each model simulates irradiation of a fissile sample, from which \gls{DNP} group parameters can be calculated.
Two types of irradiations are considered: saturation irradiations, in which long-lived \glspl{DNP} dominate, and pulse irradiations, in which short-lived \glspl{DNP} dominate.

The preliminary results cover impactful \glspl{DNP}, \gls{DNP} group parameter sensitivity studies, and 0D scaled model verification.
The results show the total delayed neutron yield and average half-life values are within uncertainty bounds between the raw \gls{DNP} data and the \gls{DNP} group parameters.
The match between the raw data and the group parameters shows the effect of each \gls{DNP} is captured in the \gls{DNP} group parameters.
Sensitivity studies reveal each \gls{DNP} affects multiple \gls{DNP} group parameters, not only the group with the most similar half-life.
Verification with the literature shows the 0D scaled model works well for longer-lived groups but poorly for shorter-lived groups.
The 0D scaled model only simulates saturation irradiations, which provide less data for shorter-lived \glspl{DNP} than pulse irradiations.

The proposed work starts with 0D flow model development, followed by a verification stage to ensure improved \gls{DNP} group parameters are correct for a static case.
The work then validates the model with a transient simulation of the \gls{MSRE} using experimental data to ensure improved \gls{DNP} group parameters are physically accurate.
The improved group spectra are then included in \gls{DNP} group parameter formulation and are verified for a static case.
Finally, the work simulates \gls{MSRE} transients for various safety, security, and safeguards scenarios using improved and pre-existing \gls{DNP} group parameters to study their impact.

This proposed work forms the critical path for a full analysis, though Section \ref{sec:future-work} presents non-critical work.
Non-critical work includes spatially resolved model development, uncertainty tracking of \gls{MSR} physical phenomena in \gls{DNP} group parameters, and investigation of the effects of varying the incident neutron energy spectrum on \gls{DNP} group parameters.